\section{Ejemplo: Problema P0}
El Problema P0 es un problema de juguete que ilustra, con un enunciado sencillo, cómo se espera que se lean los datos de entrada y se escriban las respuestas.

\subsection{Enunciado}
\textbf{Entrada:} La primera línea contiene la cantidad de casos de prueba $T$. Cada una de las siguientes $T$ líneas contiene una lista de $n$ números enteros ($0 < n \leq 1000$), separados por espacio.

\textbf{Salida:} Por cada caso de prueba se espera una línea de la forma \texttt{cp sp}, donde \texttt{cp} es la cantidad de enteros pares leídos y \texttt{sp} es la suma de esos enteros pares.

\subsection{Caso de ejemplo}
\texttt{P0.in} (entrada):
\begin{pseudocode}
6
1 2 3 4 5 6 7 8 9 10
1 -2 3 -4 5 -6 7 -8 9 -10
2 4 6 8
1 3 5 7
1 1 1 1 1
-2 2 -2 2 -2 2
\end{pseudocode}

\texttt{P0.out} (salida esperada):
\begin{pseudocode}
5 30
5 -30
4 20
0 0
0 0
6 0
\end{pseudocode}

\subsection{Ejecución}
Con los casos de prueba anteriores guardados en \texttt{P0.in}:

\textbf{Python}
\begin{pseudocode}
python ProblemaP0.py < P0.in > P0.out
\end{pseudocode}

\textbf{Java}
\begin{pseudocode}
java ProblemaP0 < P0.in > P0.out
\end{pseudocode}

\textbf{C++}
\begin{pseudocode}
./ProblemaP0 < P0.in > P0.out
\end{pseudocode}

\textbf{JavaScript}
\begin{pseudocode}
node ProblemaP0.js < P0.in > P0.out
\end{pseudocode}
